% !TEX root = ../main.tex
\lecture{6}{2026-10-19}{Grundlagen Reelle Funkitionen}

\subsection{Verknüpfung reeler Funktionen}

Operationen zwischen Abbildungsvorschriften, zB. (f+g)(x) + g(x)

\subsubsection{Max-Funktion}

Bei der max zweier Funktionen, wählt man denjenigen Wert an der gleichen stelle an beide Funktionen, der am grössten ist.


\subsection{Komposition mit affin-linearen Funktionen}

$f(x) = x - 0.001x^3$

$g(x) = 2x + 10 $4

f(g(x)) => Faktor 2x von g(x) führt dazu, dass Eingabe Werte von f(x) 2 mal schneller denselben y erreichen, also findet eine \textit{stauchung} um faktor 1/2 statt. wenn bei x = 10 den wert 30 erreicht wurde, kann jetzt mit der Komposition die x = 5 , y= 30 schon erreichen. 10 schiebt nach links aus selben grund, da x = 0 bereits den ursprunglichen wert von x = 10 erreichen kann


g(f(x)) => 2f(x) + 10, eifacher streckungsfaktor 2 und jeder Wert um 10 erhöht, also Graph nach oben verschoben.

\subsection{Grenzwerte}
\begin{definition}
    [Grenzwert einer Funktion] ähnlich wie Def. für Folgen.
    
       $ \lim_{x \rightarrow + \inf} f(x) = a$ \text{, wenn es zu jedem $\epsilon > 0$ ein $m \in \mathbb{R}$ gibt sodass für alle $x > m \exists |f(x) - a| < \epsilon$}
\end{definition}

\subsubsection{Vergleichssatz}
ähnlich wie Quetschsatz, können Grenzwerte der Terme einer gegebenen Funktion verglichen werden

\subsubsection{Rechts und Linksseitige Grenzwert}

Epsilon delta definition: 

$\lim_{x \rightarrow 10^+}$ bedeutet x nähert sich 10 an in dem es aus einer Zahl grösser immer kleiner wird. Es gibt für jedes epsilon einen delta > 0, worin alle werte grösser als 10 und weniger als delta, weniger von 10 entfernt sind als epsilon. Gleiches gilt für linksseitig.

\subsubsection{Grenzwert einer stelle}

Wenn links und rechtsseitige Grenzwerte gleich sind gibt es an der stelle x einen \textit{Grenzwert an der Stelle}, auch wenn der Wert des Grenzwerts nicht in der Definitionsmenge akzeptabel ist.
Bei unterschiede von links und rechts:
\begin{itemize}
    \item Asymptote (zb $1/x^2$ bei x = 0)
    \item Nicht stetig oder Unterbruch (Aufrundungsfunktion)
\end{itemize}

\begin{definition}
    [Stetigkeit eines Funktionpunkts] wenn in $x_0$ der $lim{x\rightarrow x_0} = f(x_0)$ oder wenn an der Stelle x0 der Grenzwert existiert und gleich dem Funktionwert an der Stelle ist.
\end{definition}

\begin{definition}
    [Stückweise stetige Funktion] wenn endliche Anzahl an Unstetigkeiten vorhanden 
\end{definition}